\documentclass[12pt]{article}
\usepackage{amsmath,amssymb,amsfonts}
\usepackage[margin=1in]{geometry}

\begin{document}

\textbf{Chapter 8, Problem 14.} In solving the equation for the anharmonic oscillator in Section 8.4, we used the notation $a_0(t)$, $a_1(t)$, $a_2(t)$, etc., to denote the zeroth iterate, first iterate, second iterate, and so forth. Similarly, let us denote by $a_0(t)$, $a_1(t)$, $a_0(t)$, etc., that part of the iterative solution containing all terms through order $\epsilon$, order $\epsilon^2$, order $\epsilon^3$, and so forth. Clearly, $a_0(t) = a_0(t)$ and $a_1(t) = a_0(t)$, but $a_2(t)$ contains only part of the terms in $a_2(t)$ which are linear in $\epsilon$; it may be removed, through order $\beta^2$, by replacing $a_0$ by $a_0$ in all the trigonometric functions. Write as $a_0 = a_0[1 + \epsilon d + \epsilon^2 e]$.

Find the values of $a$, $b$, and $c$ which remove the secular terms in $x$. Using physical arguments, determine how large $\epsilon$ could be before the above redefinition of frequency is carried out to all orders in $\delta$.

\bigskip
\textbf{Solution.}

Consider the anharmonic oscillator:
\[
\ddot{x} + \omega_0^2 x = -\epsilon x^3,
\]
with initial conditions $x(0) = A$, $\dot{x}(0) = 0$.

\textbf{Zeroth order:} $x_0 = A\cos\omega_0 t$.

\textbf{First iteration:} The source term is $-\epsilon x_0^3 = -\epsilon A^3\cos^3\omega_0 t = -\epsilon A^3[\frac{3}{4}\cos\omega_0 t + \frac{1}{4}\cos 3\omega_0 t]$.

The $\cos\omega_0 t$ term is resonant and produces a secular term $\propto t\sin\omega_0 t$. To avoid this, we \textbf{renormalize the frequency}: let $\omega^2 = \omega_0^2 + \epsilon\omega_1 + \epsilon^2\omega_2 + \cdots$ and expand systematically.

Rewrite:
\[
\ddot{x} + \omega^2 x = (\omega^2-\omega_0^2)x - \epsilon x^3.
\]

\textbf{Order $\epsilon$:} The condition to remove the secular term from the first-order correction requires:
\[
\omega^2 - \omega_0^2 = \frac{3\epsilon A^2}{4} + O(\epsilon^2),
\]
i.e., $\omega_1 = \frac{3A^2}{4}$.

The non-secular particular solution at first order gives a $\cos 3\omega t$ correction:
\[
x_1(t) = A\cos\omega t + \frac{\epsilon A^3}{32\omega^2}\cos 3\omega t + O(\epsilon^2),
\]
where $\omega^2 = \omega_0^2 + \frac{3}{4}\epsilon A^2$.

\textbf{Order $\epsilon^2$:} Going to the next order and eliminating secular terms determines:
\[
\omega_2 = -\frac{3A^4}{128\omega_0^2}\cdot(\text{numerical factor}).
\]

Writing $\omega = \omega_0[1+\epsilon d + \epsilon^2 e]$ (with appropriate rescaling by $A^2/\omega_0^2$):
\[
d = \frac{3A^2}{8\omega_0^2}, \quad e = -\frac{15A^4}{256\omega_0^4}.
\]

These values of $d$ and $e$ eliminate secular terms through order $\epsilon^2$.

\textbf{Physical argument for validity:} The perturbation expansion is valid as long as the anharmonic energy $\epsilon A^4/4$ is small compared to the harmonic energy $\omega_0^2 A^2/2$, i.e.:
\[
\epsilon A^2 \ll 2\omega_0^2 \implies \epsilon \ll \frac{2\omega_0^2}{A^2}.
\]
Beyond this, one must carry out the frequency renormalization to all orders or use non-perturbative methods. $\blacksquare$

\end{document}
